\begin{textsample}{POS Dim 6 – sp – Score 8.00 – sp/sp\_ef\_ciencias\_da\_natureza\_1.txt}  \label{ex:f6_pos_085}
Área de Ciências da Natureza O conhecimento científico e tecnológico intervém no modo de vida e na forma como a sociedade se organiza contemporaneamente.
Isto exige investir na formação de um sujeito transformador do seu meio, que reflita, proponha, argumente e aja com base em fundamentos científicos e tecnológicos, de modo intencional e consciente, em todos os âmbitos da vida humana.
Portanto, ao longo do Ensino Fundamental, a área de Ciências da Natureza tem um compromisso com o desenvolvimento do Letramento Científico, que envolve a capacidade de compreender e interpretar o mundo ( natural, social e tecnológico ), mas também de transformá -lo com base nos aportes teóricos e processuais das ciências.
Nessa perspectiva, por meio de um olhar articulado de diversos campos do saber, a área pretende assegurar aos estudantes o acesso à diversidade de conhecimentos científicos produzidos ao longo da história, bem como a aproximação gradativa aos principais processos, práticas e procedimentos da Investigação Científica.
No Currículo Paulista, as habilidades da área estão relacionadas de modo a construir e consolidar conhecimentos, desde a Educação Infantil, passando pelo Ensino Fundamental, até o Ensino Médio, com vistas ao Letramento Científico, na perspectiva anteriormente explicitada.
Para o desenvolvimento dessas habilidades, alguns princípios são fundamentais.
O primeiro deles ressalta a necessidade de considerar o contexto das aprendizagens da área.
A construção e a consolidação do conhecimento científico devem, sempre que possível, estabelecer relação com as experiências vivenciadas pelos estudantes nos diversos espaços que constituem sua vida e seu cotidiano.
Isso implica a necessidade de fundamentar e correlacionar os conhecimentos construídos ao conhecimento científico, de modo que os estudantes possam constituir estruturas explicativas importantes para significar aquilo que aprendem e criar condições para que possam validar o conhecimento científico envolvido em sua experiência escolar.
É necessário, ainda que progressivamente, que possam apropriar -se da Linguagem Científica.
Na área de Ciências da Natureza, valorizar a experiência de aprendizagem de cada estudante implica conceber o ensino por meio da investigação.
Trata -se de desenvolver as aprendizagens, recorrendo aos procedimentos de investigação em todos os anos da Educação Básica, sendo este outro princípio orientador da área.
A investigação pressupõe a observação, a análise de evidências e proposição de hipóteses na definição de um \textbf{problema}, a experimentação, a construção de modelos, entre outros processos e métodos.
Nesse exercício investigativo podem ser desenvolvidos o pensamento crítico, a criatividade, a responsabilidade e a autonomia, bem como aprofundar as relações interpessoais.
O estudante experimenta, pesquisa, levanta hipóteses científicas, \textbf{testa} essas hipóteses, aprende a problematizar, argumentar e olhar criticamente para todos os fenômenos ( naturais ou sociais ), para si mesmo e para o outro.
Cabe ressaltar que, segundo a Base Nacional Comum Curricular ( BNCC ), adotar os procedimentos de investigação não significa realizar atividades seguindo, necessariamente, um conjunto de etapas predefinidas, tampouco restringe -se à mera manipulação de objetos ou realização de experimentos em laboratório.
É imprescindível que os estudantes sejam progressivamente estimulados e apoiados na proposição de situações a serem investigadas, no planejamento e na realização colaborativa de atividades investigativas, bem como no compartilhamento e na comunicação dos resultados dessas investigações.
Além disso, é desejável que aprendam a valorizar erros e acertos desses processos, assim como possam propor intervenções orientadas pelos resultados obtidos, com foco na melhoria da qualidade de vida individual e coletiva, da saúde, da \textbf{sustentabilidade} e/ou na resolução de \textbf{problemas} cotidianos.
Dessa maneira, os estudantes podem consolidar e ampliar as concepções sobre fatos e fenômenos da natureza de modo a compreender melhor o ambiente, numa perspectiva ecológica e social, considerando os aspectos econômicos e políticos que se articulam e se manifestam no âmbito local e global.
Da mesma forma, podem avaliar os impactos \textbf{ambientais} nas áreas do trabalho, da \textbf{tecnologia}, da produção de energia, da \textbf{sustentabilidade}, da urbanização e do campo.
Sendo assim, em relação aos procedimentos de investigação, o ensino de Ciências deve promover situações nas quais os estudantes possam : Procedimentos de investigação | Procedimentos | | |---|---| | Definição de \textbf{Problemas} | • Observar o mundo a sua volta e fazer perguntas ; • Analisar demandas, delinear \textbf{problemas} e planejar investigações ; • Propor hipóteses. | | Levantamento, Análise e Representação | • Planejar e realizar atividades de campo ( experimentos, observações, leituras, visitas, ambientes virtuais etc. ) ; • Desenvolver e utilizar ferramentas, inclusive \textbf{digitais}, para coleta, análise e representação de dados ( imagens, esquemas, tabelas, gráficos, quadros, diagramas, mapas, modelos, representações de sistemas, fluxogramas, mapas conceituais, simulações, \textbf{aplicativos} etc. ) ; • Avaliar a informação ( validade, coerência e adequação ao \textbf{problema} formulado ) ; • Elaborar explicações e/ou modelos ; • Associar explicações e/ou modelos à evolução histórica dos conhecimentos científicos envolvidos ; • Selecionar e construir \textbf{argumentos} com base em evidências, modelos e/ou conhecimentos científicos ; • Aprimorar seus saberes e incorporar, gradualmente, e de modo significativo, o conhecimento científico ; • Desenvolver soluções para \textbf{problemas} cotidianos usando diferentes ferramentas, inclusive \textbf{digitais}. | | Comunicação | • Organizar e/ou extrapolar conclusões ; • Relatar informações de forma oral, escrita ou multimodal ; • Apresentar, de forma sistemática, dados e resultados de investigações ; • Participar de discussões de caráter científico com colegas, professores, familiares e comunidade em geral ; • Considerar contra-argumentos para rever processos investigativos e conclusões. | | Intervenção | • Implementar soluções e avaliar sua eficácia para resolver \textbf{problemas} cotidianos ; • Desenvolver ações de intervenção para melhorar a qualidade de vida individual, coletiva e socioambiental. |

% matched lemmas: ambiental, aplicativo, argumento, digital, problema, sustentabilidade, tecnologia, testar
\end{textsample}
